\chapter{GIỚI THIỆU}

\section{Giới thiệu đề tài}

Trong bối cảnh trí tuệ nhân tạo (Artificial Intelligence - AI) phát triển mạnh mẽ, các mô hình ngôn ngữ lớn (Large Language Models - LLM) đang được ứng dụng rộng rãi trong nhiều lĩnh vực như giáo dục, y tế, tài chính và dịch vụ khách hàng. Đặc biệt, các hệ thống Chatbot thông minh đã trở thành công cụ hỗ trợ hiệu quả trong việc tìm kiếm, tổng hợp và trả lời thông tin.

Trong lĩnh vực giáo dục, sinh viên thường gặp khó khăn khi phải tìm kiếm thông tin từ nhiều nguồn tài liệu khác nhau như giáo trình, slide bài giảng, tài liệu PDF hoặc tài liệu tham khảo. Việc đọc và tổng hợp lượng lớn dữ liệu mất nhiều thời gian và đôi khi không mang lại kết quả chính xác.

Xuất phát từ nhu cầu đó, nhóm thực hiện đề tài:

\begin{center}
	\textbf{\Large XÂY DỰNG HỆ THỐNG CHATBOT HỖ TRỢ HỌC TẬP ỨNG DỤNG RETRIEVAL-AUGMENTED GENERATION (RAG)}
\end{center}

Hệ thống được xây dựng với mục tiêu hỗ trợ sinh viên đặt câu hỏi trực tiếp về nội dung môn học. Chatbot sẽ truy xuất thông tin từ tài liệu đã được tải lên, kết hợp cùng mô hình ngôn ngữ lớn để tạo ra câu trả lời có độ chính xác cao và sát với tài liệu gốc.

Bên cạnh đó, đề tài còn nghiên cứu, đánh giá và so sánh hiệu quả giữa hai phương pháp phổ biến hiện nay là Retrieval-Augmented Generation (RAG) và Fine-Tuning nhằm xác định giải pháp phù hợp cho bài toán hỏi đáp tài liệu học tập.

\section{Lý do chọn đề tài}

Hiện nay, phần lớn sinh viên sử dụng các công cụ AI như ChatGPT hoặc Gemini để hỗ trợ học tập. Tuy nhiên, các mô hình này chủ yếu dựa trên dữ liệu được huấn luyện trước và không thể truy cập trực tiếp vào tài liệu nội bộ của từng môn học nếu không được cung cấp.

Điều này dẫn đến một số hạn chế:

\begin{itemize}
	\item Câu trả lời có thể không đúng với giáo trình của môn học.
	\item Khó trích dẫn nguồn tài liệu.
	\item Không thể cập nhật nhanh khi tài liệu môn học thay đổi.
	\item Độ chính xác giảm đối với các kiến thức chuyên ngành.
\end{itemize}

Trong khi đó, kỹ thuật Retrieval-Augmented Generation (RAG) cho phép hệ thống truy xuất trực tiếp thông tin từ các tài liệu do người dùng cung cấp trước khi sinh câu trả lời. Đồng thời, Fine-Tuning giúp mô hình học sâu hơn trên một tập dữ liệu chuyên biệt để nâng cao chất lượng phản hồi.

Vì vậy, nhóm lựa chọn nghiên cứu và xây dựng hệ thống Chatbot kết hợp RAG và Fine-Tuning nhằm hỗ trợ học tập, đồng thời đánh giá ưu nhược điểm của từng phương pháp trong thực tế.

\section{Mục tiêu đề tài}

Đề tài hướng đến các mục tiêu sau:

\begin{itemize}
	
	\item Xây dựng hệ thống Chatbot hỗ trợ trả lời câu hỏi liên quan đến tài liệu môn học.
	
	\item Cho phép người quản trị tải lên và quản lý tài liệu PDF phục vụ việc hỏi đáp.
	
	\item Xây dựng quy trình xử lý tài liệu bao gồm tách văn bản, chia nhỏ dữ liệu (Chunking), tạo Embedding và lưu trữ trong cơ sở dữ liệu vector ChromaDB.
	
	\item Tích hợp mô hình ngôn ngữ lớn thông qua FastAPI để sinh câu trả lời.
	
	\item Xây dựng giao diện Web bằng React giúp người dùng dễ dàng tương tác với hệ thống.
	
	\item So sánh chất lượng giữa mô hình RAG và Fine-Tuning dựa trên các tiêu chí như độ chính xác, tốc độ phản hồi và khả năng bám sát tài liệu.
	
	\item Đánh giá khả năng ứng dụng của hệ thống trong môi trường học tập thực tế.
	
\end{itemize}

\section{Phạm vi nghiên cứu}

Đề tài tập trung nghiên cứu và triển khai các nội dung sau:

\begin{itemize}
	
	\item Xây dựng Chatbot phục vụ lĩnh vực học tập.
	
	\item Hỗ trợ truy vấn thông tin từ tài liệu PDF.
	
	\item Triển khai kỹ thuật RAG kết hợp ChromaDB.
	
	\item Nghiên cứu Fine-Tuning trên tập dữ liệu môn học.
	
	\item Phát triển Backend bằng FastAPI.
	
	\item Phát triển Frontend bằng ReactJS.
	
	\item Đánh giá hiệu năng và độ chính xác giữa RAG và Fine-Tuning.
	
\end{itemize}

Đề tài chưa nghiên cứu các nội dung như nhận dạng giọng nói, xử lý hình ảnh, đa ngôn ngữ hoặc triển khai trên môi trường Cloud quy mô lớn.

\section{Đối tượng sử dụng}

Hệ thống được xây dựng dành cho hai nhóm người dùng chính:

\subsection*{Sinh viên}

\begin{itemize}
	\item Đặt câu hỏi về nội dung môn học.
	\item Tra cứu thông tin nhanh.
	\item So sánh kết quả giữa RAG và Fine-Tuning.
	\item Xem báo cáo đánh giá chất lượng câu trả lời.
\end{itemize}

\subsection*{Quản trị viên}

\begin{itemize}
	\item Quản lý tài liệu môn học.
	\item Cập nhật nguồn dữ liệu.
	\item Theo dõi trạng thái hệ thống.
	\item Xem báo cáo thống kê.
	\item Quản lý quá trình Fine-Tuning.
\end{itemize}

\section{Phương pháp nghiên cứu}

Đề tài được thực hiện theo các bước:

\begin{enumerate}
	
	\item Khảo sát các mô hình Chatbot hiện nay.
	
	\item Nghiên cứu kiến trúc RAG.
	
	\item Nghiên cứu phương pháp Fine-Tuning.
	
	\item Thiết kế kiến trúc hệ thống.
	
	\item Xây dựng Backend bằng FastAPI.
	
	\item Xây dựng Frontend bằng React.
	
	\item Tích hợp ChromaDB và Embedding.
	
	\item Kiểm thử chức năng.
	
	\item Đánh giá kết quả.
	
\end{enumerate}

\section{Cấu trúc báo cáo}

Báo cáo gồm các chương sau:

\begin{itemize}
	
	\item Chương 1: Giới thiệu.
	
	\item Chương 2: Cơ sở lý thuyết và công nghệ sử dụng.
	
	\item Chương 3: Đặc tả yêu cầu phần mềm (SRS).
	
	\item Chương 4: Thiết kế hệ thống.
	
	\item Chương 5: Triển khai hệ thống.
	
	\item Chương 6: Kiểm thử và đánh giá.
	
	\item Chương 7: Quản lý dự án bằng Jira và GitHub.
	
	\item Chương 8: Kết luận và hướng phát triển.
	
\end{itemize}

